\documentclass[11pt,a4paper]{article}
\usepackage[utf8]{inputenc}
\usepackage[T1]{fontenc}
\usepackage[margin=1in]{geometry}
\usepackage{xcolor}
\usepackage{hyperref}

\definecolor{primarycolor}{RGB}{33,37,41}
\definecolor{secondarycolor}{RGB}{0,102,204}
\hypersetup{colorlinks=true,linkcolor=secondarycolor,urlcolor=secondarycolor}
\pagestyle{empty}
\setlength{\parskip}{0.9em}
\setlength{\parindent}{0pt}

\begin{document}

Dear Hiring Manager,

I am applying for the Cloud Infrastructure / DevOps Contractor role supporting Corvale Mobility’s fleet connectivity platform. In my current contract as a Cloud Engineer, I build Azure infrastructure with Terraform, operate AKS, and create deployment pipelines that help move services into production reliably.

At Havenstad Logistics, I designed a hub-and-spoke Azure network spanning three subscriptions, with Azure Firewall, private DNS, private endpoints, and peering. I also added Terratest module tests and nightly drift detection. For internal applications, I replaced manual releases with Azure DevOps pipelines for 15 .NET and Node apps, raising deployment frequency from roughly monthly to several times per week for most apps. This work has given me practical experience with infrastructure design, networking, automation, and supporting application delivery.

Earlier, I planned and carried out a migration of 60 servers from VMware to Azure IaaS, completing the cut-over over six weekends without unplanned downtime. My systems administration background also includes Windows and Linux environments, while my freelance work included PowerShell-based Microsoft 365 security baselines. Corvale’s project to connect a new generation of trains with cloud services is a compelling opportunity to apply this infrastructure experience to a platform being shaped from the start and taken into production.

I would welcome the opportunity to discuss how my Azure, Terraform, networking, and delivery-pipeline experience could contribute to Corvale’s fleet connectivity work. Thank you for your consideration.

Kind regards,
Samuel Achterberg-Osei

\end{document}